\section{Data Architecture}
\label{sec:data-architecture}

\subsection{Design Principle: Separation of Visual and Routing Layers}

The system data model enforces a strict separation between two concerns:

\begin{enumerate}
  \item \textbf{Visual Layer} --- static image assets (floor plan PNGs or SVGs) used purely for rendering in the browser canvas. These carry no routing logic.
  \item \textbf{Routing Layer} --- a graph-based data structure of nodes and edges that encodes all navigable paths, independent of the visual representation.
\end{enumerate}

This decoupling follows the hybrid spatial model principle established by Lorenz et al.\ \cite{Lorenz2006}, who demonstrated that indoor environments benefit from maintaining separate geometric and topological representations. Liu and Zlatanova \cite{Liu2012} further validated this approach with a semantic data model that cleanly isolates navigation networks from spatial geometry.

The advantage for a web application is clear: the floor plan image can be replaced or updated without modifying routing data, and the routing graph can be edited without regenerating visual assets.

\subsection{Floor Schema}

Each floor in the building is represented as a discrete record linking metadata to its visual asset and coordinate calibration:

\begin{lstlisting}[language=json, caption={Floor data structure}]
{
  "floors": [
    {
      "floor_id": "F1",
      "label": "Ground Floor",
      "elevation": 0.0,
      "image_asset": "assets/floor1.svg",
      "image_width_px": 1200,
      "image_height_px": 800,
      "scale_factor": 0.05
    },
    {
      "floor_id": "F2",
      "label": "Second Floor",
      "elevation": 4.0,
      "image_asset": "assets/floor2.svg",
      "image_width_px": 1200,
      "image_height_px": 800,
      "scale_factor": 0.05
    },
    {
      "floor_id": "F3",
      "label": "Third Floor",
      "elevation": 8.0,
      "image_asset": "assets/floor3.svg",
      "image_width_px": 1200,
      "image_height_px": 800,
      "scale_factor": 0.05
    }
  ]
}
\end{lstlisting}

\texttt{scale\_factor} is defined as meters-per-pixel and is used by the distance approximation module (Section~\ref{sec:distance-approximation}). \texttt{elevation} records the real-world height in meters above ground datum for each floor, enabling vertical distance cost calculation.

\subsection{Node Schema}

Nodes represent discrete navigable points within the building. Following the topological routing-graph approach described by Willemsen et al.\ \cite{Willemsen2015}, each node carries:

\begin{lstlisting}[language=json, caption={Node data structure}]
{
  "nodes": [
    {
      "node_id": "N1_ENTRANCE",
      "floor_id": "F1",
      "x": 100,
      "y": 750,
      "type": "entrance",
      "label": "Main Entrance"
    },
    {
      "node_id": "N1_HALL_A",
      "floor_id": "F1",
      "x": 300,
      "y": 500,
      "type": "corridor",
      "label": null
    },
    {
      "node_id": "N1_ROOM101",
      "floor_id": "F1",
      "x": 500,
      "y": 400,
      "type": "room",
      "label": "Room 101 - Registrar"
    },
    {
      "node_id": "N1_STAIR_A",
      "floor_id": "F1",
      "x": 600,
      "y": 500,
      "type": "vertical",
      "label": "Stairwell A (F1)"
    },
    {
      "node_id": "N2_STAIR_A",
      "floor_id": "F2",
      "x": 600,
      "y": 500,
      "type": "vertical",
      "label": "Stairwell A (F2)"
    },
    {
      "node_id": "N2_ROOM201",
      "floor_id": "F2",
      "x": 400,
      "y": 300,
      "type": "room",
      "label": "Room 201 - Library"
    }
  ]
}
\end{lstlisting}

Node types are enumerated as:

\begin{table}[H]
\centering
\caption{Node type classification}
\begin{tabularx}{0.85\textwidth}{lX}
\toprule
\textbf{Type} & \textbf{Description} \\
\midrule
\texttt{room} & A navigable destination (selectable by the user) \\
\texttt{corridor} & An intermediate routing waypoint (not user-selectable) \\
\texttt{vertical} & A floor-transition point: stairwell, elevator, or ramp \\
\texttt{entrance} & A building entry/exit point \\
\bottomrule
\end{tabularx}
\end{table}

The \texttt{x} and \texttt{y} coordinates are pixel positions relative to the top-left origin of the corresponding floor image. This aligns with the Canvas 2D API coordinate system used in the browser rendering layer.

\subsection{Edge Schema}

Edges encode navigable connections between nodes. Zhang and Ye \cite{Zhang2017} demonstrated that graph-based indoor models require explicit encoding of both horizontal and vertical transitions as weighted edges:

\begin{lstlisting}[language=json, caption={Edge data structure}]
{
  "edges": [
    {
      "edge_id": "E001",
      "from_node": "N1_ENTRANCE",
      "to_node": "N1_HALL_A",
      "weight": 12.5,
      "is_vertical": false,
      "is_accessible": true
    },
    {
      "edge_id": "E002",
      "from_node": "N1_HALL_A",
      "to_node": "N1_ROOM101",
      "weight": 8.3,
      "is_vertical": false,
      "is_accessible": true
    },
    {
      "edge_id": "E003",
      "from_node": "N1_HALL_A",
      "to_node": "N1_STAIR_A",
      "weight": 15.0,
      "is_vertical": false,
      "is_accessible": true
    },
    {
      "edge_id": "E004",
      "from_node": "N1_STAIR_A",
      "to_node": "N2_STAIR_A",
      "weight": 6.0,
      "is_vertical": true,
      "is_accessible": false
    },
    {
      "edge_id": "E005",
      "from_node": "N2_STAIR_A",
      "to_node": "N2_ROOM201",
      "weight": 10.2,
      "is_vertical": false,
      "is_accessible": true
    }
  ]
}
\end{lstlisting}

Key design decisions for edges:

\begin{itemize}
  \item \textbf{Weight}: Pre-computed distance in meters. For horizontal edges, this is derived from the Euclidean pixel distance multiplied by the floor's \texttt{scale\_factor}. For vertical edges, the weight includes an estimated traversal cost (e.g., stair-climbing time equivalent).
  \item \textbf{Bidirectional}: All edges are treated as bidirectional by default. The routing engine traverses them in both directions unless explicitly restricted.
  \item \textbf{Vertical flag}: Marks edges that cross floors. The routing engine uses this flag to trigger floor-transition logic in the UI (Section~\ref{sec:routing-engine}).
  \item \textbf{Accessibility flag}: Enables filtering for wheelchair-accessible routes by excluding stair-only vertical edges.
\end{itemize}

\subsection{Graph Summary for Target Building}

For the target 3-story, 5-rooms-per-floor building, the expected graph dimensions are:

\begin{table}[H]
\centering
\caption{Estimated graph size}
\begin{tabular}{lrl}
\toprule
\textbf{Component} & \textbf{Count} & \textbf{Rationale} \\
\midrule
Room nodes     & 15 & 5 rooms $\times$ 3 floors \\
Corridor nodes & $\sim$20 & Hallway waypoints per floor \\
Vertical nodes & $\sim$6 & 2 stairwells $\times$ 3 floors \\
Entrance nodes & 1--2 & Building entry points \\
\midrule
Total nodes    & $\sim$43 & \\
Total edges    & $\sim$60 & \\
\bottomrule
\end{tabular}
\end{table}

This is a trivially small graph that fits entirely in a static JSON file served to the browser---no database or server-side computation is required. The entire data payload is estimated at under 5\,KB uncompressed.
